\documentclass[11pt,a4paper]{article}
\usepackage{fontspec,amsmath,amssymb,amsthm,microtype}
\setmainfont{lmroman10-regular.otf}[BoldFont=lmroman10-bold.otf,ItalicFont=lmroman10-italic.otf,BoldItalicFont=lmroman10-bolditalic.otf]
\usepackage[margin=25mm]{geometry}
\usepackage[unicode,hidelinks]{hyperref}
\hypersetup{pdfauthor={LangevinLean Project}}
\setlength{\emergencystretch}{3em}
\setlength{\parskip}{0.4em}
\setlength{\parindent}{0pt}
\allowdisplaybreaks[2]
\raggedbottom
\title{\texorpdfstring{Scaled coordinates for the parameter and particles}{Scaled coordinates for the parameter and particles}}
\author{LangevinLean Project}
\date{}
\begin{document}\maketitle
Draft mathematical exposition; these notes may contain mistakes.


Use orthogonal coordinates for the common parameter and the scaled latent particles. Their norm satisfies \(\|z\|^2=\|\theta\|^2+N^{-1}\sum_i\|x_i\|^2\). The average over physical particle energies is therefore the squared Euclidean norm of the latent blocks. This relation determines the particle factors in the convexity, gradient, and Euler-update formulas.



\section{\texorpdfstring{The coordinate spaces and their dimensions}{The coordinate spaces and their dimensions}}


A scaled particle state belongs to \(\mathbb R^{d_\theta}\times(\mathbb R^{d_x})^N\); a parameter–particle pair belongs to \(\mathbb R^{d_\theta}\times\mathbb R^{d_x}\). The parameter block and the latent blocks are mutually orthogonal, so their dimensions add:

\[\dim(\mathrm{State})=d_\theta+Nd_x,\qquad \dim(\mathrm{Joint})=d_\theta+d_x.\]

The block map returns \(z_i\), and its corresponding physical pair is \(A_i z=\operatorname{joint}_i(z)=(\theta,\sqrt N z_i)\). The potential \(V_N\) is the sum of the pair potentials at these physical coordinates. Thus \(\mathbb R^{d_\theta}\times(\mathbb R^{d_x})^N\) contains one common parameter and one scaled block for each particle. Changing the parameter changes every pair, whereas changing a single latent block changes only that pair. This distinction determines which gradient coordinates are summed and which acquire a factor of \(N\).



\section{\texorpdfstring{A sum of squared pair norms}{A sum of squared pair norms}}


The extraction map \(A_i\) is linear because its parameter and latent coordinates are linear functions of the state. In the sum of the extracted squared norms, each parameter coordinate occurs \(N\) times and each latent coordinate is multiplied once by \(\sqrt N\). Expanding the norms therefore gives

\[\sum_i\|A_i z\|^2=N\|\theta\|^2+N\sum_i\|z_i\|^2=N\|z\|^2.\]

The equality uses \((\sqrt N)^2=N\) and remains valid when \(N=0\), since both finite sums then have the required zero contributions. Linearity gives the corresponding identity for two states, \(\sum_i\|A_i(z-w)\|^2=N\|z-w\|^2\). Applying this sum of squares in the convexity and gradient estimates retains the precise particle factor; estimating the extraction maps separately would give a less precise bound.



\section{\texorpdfstring{Strong convexity of the summed pair potential}{Strong convexity of the summed pair potential}}


Assume that \(U\) is \(m\)-strongly convex. For weights \(a,b\ge0\) satisfying \(a+b=1\), its defining inequality at each pair is

\[U(aA_ix+bA_iy)\le aU(A_ix)+bU(A_iy)-\frac{mab}{2}\|A_i(x-y)\|^2.\]

Sum over \(i\). Linearity identifies the left-hand arguments with \(A_i(ax+by)\), and the squared-norm identity applied to \(x-y\) identifies the sum of quadratic terms. Consequently \(V_N\) is strongly convex with constant \(Nm\). The map \(A_i\) is continuous because it is linear between finite-dimensional spaces. If \(U\) is differentiable, the chain rule and the finite sum imply differentiability of the sum of \(U\circ A_i\).



More explicitly, the summed quadratic correction is \((mab/2)\sum_i\|A_i(x-y)\|^2\), which equals \((Nmab/2)\|x-y\|^2\). The curvature factor is therefore exactly \(N\) times the pair curvature. The argument includes weights zero and one, for which the correction vanishes, and requires no differentiability for the convexity conclusion.



\section{\texorpdfstring{The gradient chain rule and the lifted Lipschitz constant}{The gradient chain rule and the lifted Lipschitz constant}}


Assume that the pair potential is differentiable. In a direction \(w\), the chain rule gives

\[\langle\nabla V_N(z),w\rangle=\sum_i\langle\nabla U(A_i z),A_iw\rangle.\]

Cauchy–Schwarz in the particle index and the squared-norm identity give \(\sum_i\|A_iv\|\|A_iw\|\le N\|v\|\|w\|\). If \(\nabla U\) is \(L\)-Lipschitz, subtraction of the directional formulas at two states yields

\[\langle\nabla V_N(x)-\nabla V_N(y),w\rangle\le NL\|x-y\|\|w\|.\]

Set \(w=\nabla V_N(x)-\nabla V_N(y)\). A zero gradient difference already satisfies the desired estimate; otherwise division by its positive norm shows that \(\nabla V_N\) is \(NL\)-Lipschitz. Applying the same argument to the absolute value gives \(|\langle\nabla V_N(x)-\nabla V_N(y),w\rangle|\le NL\|x-y\|\|w\|\), which is also the dual-norm form of the estimate. Both curvature and gradient regularity are multiplied by \(N\), so the condition number is \(L/m\) whenever the curvature is positive.




\section{\texorpdfstring{The two gradient blocks}{The two gradient blocks}}


Let \(e_j\) be a parameter basis vector. Each \(A_i e_j\) is the corresponding parameter basis vector in the pair space. For the latent basis vector of particle \(i\), one has \(A_qe_{i,j}=0\) whenever \(q\ne i\), while \(A_i e_{i,j}=\sqrt N e_j\). Substitution in the directional derivative formula gives

\[\nabla_\theta V_N(z)=\sum_i\nabla_\theta U(\theta,\sqrt Nz_i),\qquad (\nabla V_N(z))_i=\sqrt N\nabla_xU(\theta,\sqrt Nz_i).\]

These are the two coordinate derivatives before multiplication by the step divided by \(N\). The parameter derivative of \(V_N\) contains every summand because the parameter is common to all pairs. For a latent block, the substitution \(x_i=\sqrt N z_i\) contributes a derivative factor \(\sqrt N\) to its single summand. Multiplication by \(\gamma/N\) then produces the parameter average and the physical-particle drift.



\section{\texorpdfstring{The parameter projection is a contraction}{The parameter projection is a contraction}}


The parameter projection \(p_\theta z=\theta\) is linear and continuous. Orthogonality of the blocks gives \(\|p_\theta z\|^2\le\|z\|^2\) because every omitted squared coordinate is nonnegative. Taking square roots gives \(\|p_\theta z\|\le\|z\|\), and the definition of the operator norm gives \(\|p_\theta\|\le1\).



For two states, linearity also gives \(\|p_\theta z-p_\theta w\|\le\|z-w\|\). Applying this inequality to the two coordinates of a coupling shows that parameter projection cannot increase its quadratic cost. The inequality remains valid after taking an \(L^2\) norm. It therefore converts a full-state transport estimate into an estimate for the parameter, to which the equilibrium error about a deterministic center can then be added.



\section{\texorpdfstring{Equivalence of the scaled and physical updates}{Equivalence of the scaled and physical updates}}


For physical step \(h\), define the scaled Euler update by

\[z'=z-\frac hN\nabla V_N(z)+\sqrt{2h/N}\,\xi.\]

The parameter coordinate is \(\theta'=\theta-(h/N)\sum_i\nabla_\theta U(\theta,x_i)+\sqrt{2h/N}\,\xi_\theta\). Suppose \(N>0\). Multiplication of the \(i\)-th latent coordinate by \(\sqrt N\) gives drift coefficient \((\sqrt N)^2h/N=h\) and noise coefficient \(\sqrt N\sqrt{2h/N}=\sqrt{2h}\). Thus the physical latent update is

\[x_i'=x_i-h\nabla_xU(\theta,x_i)+\sqrt{2h}\,\xi_i.\]

The factors arise from the pair derivatives: the latent derivative contains \(\sqrt N\), and conversion back to physical coordinates contributes another \(\sqrt N\). For step notation used in the probabilistic algorithm, \((\gamma/N)\sqrt N\sqrt N\) is therefore \(\gamma\) and the latent noise coefficient is \(\sqrt{2\gamma}\). The parameter noise coefficient remains \(\sqrt{2\gamma/N}\). Positivity of \(N\) justifies the divisions.



For \(\gamma>0\), the Gaussian parameter variance in each coordinate is \(2\gamma/N\), while each physical latent variance is \(2\gamma\). These are precisely the variances of the interacting particle updates. The coordinate equivalence is an equality of updates and does not involve a discretization approximation.




\end{document}
